
un homomorphisme de spécialisation :
@@M0@@
en définissant pour tout cycle premier @@M1@@ de @@M2@@, purement de dimension @@M3@@, identifié à un sous-schéma intègre de @@M4@@, l'élément @@M5@@ de @@M6@@ comme le cycle associé au sous-schéma fermé @@M7@@ de @@M8@@, où @@M9@@ est l'adhérence schématique de @@M10@@ dans @@M11@@ (qui est un sous-schéma fermé de @@M12@@ plat sur @@M13@@ (EGA IV 2.8), de fibre générique égale au @@M14@@ de départ, et satisfaisant @@M15@@ (EGA IV 7.1.13)). Lorsque @@M16@@ est localement équidimensionnel (i.e. toutes ses composantes irréductibles passant par un point donné sont de même dimension) on vérifie aussitôt que l'homomorphisme de spécialisation qu'on vient de définir induit des homomorphismes
@@M17@@

Supposons maintenant @@M18@@ quasi-projectif, de sorte qu'on dispose de la notion d'équivalence rationnelle des cycles sur @@M19@@ et sur @@M20@@ (cf. [5], où la théorie est faite sur un corps algébriquement clos). Il semble très plausible que l'homomorphisme (7.14.1) transforme cycles rationnellement équivalents à zéro en cycles rationnellement équivalents à zéro, et induise donc un homomorphisme correspondant sur les groupes «~de Chow~» de classes de cycles
@@M21@@
ou encore, lorsque @@M22@@ est localement équidimensionnel :
@@M23@@

On pourrait songer à le démontrer par une méthode paraphrasant celle de 7.12, mais où le recours à la résolution des singularités et la réduction
